\item \textbf{Problema de repaso.} El géiser Old Faithful en el parque nacional Yellowstone erupciona a intervalos aproximados de $1 \text{ hora}$, y la altura de la columna de agua alcanza $40.0 \text{ m}$ (figura P14.25). (a) Modele la corriente que se eleva como una serie de gotas separadas. Analice el movimiento en caída libre de una de las gotas para determinar la rapidez a la que el agua deja el suelo. (b) ¿Qué pasaría si? Modele la corriente que se eleva como un fluido ideal en un flujo de líneas de corriente. Use la ecuación de Bernoulli para determinar la rapidez del agua mientras deja el nivel del suelo. (c) ¿De qué modo se compara la respuesta del inciso (a) con la respuesta del inciso (b)? (d) ¿Cuál es la presión (arriba de la atmosférica) en la cámara subterránea caliente si su profundidad es de $175 \text{ m}$? Suponga que la cámara es grande en comparación con la boca del géiser.
